\begin{claim}{D.2.29}{Proved}{A fourth-suspension unit map and its canonical obstruction}
Let $B=K(\FF_2,2)$, $X=\Sigma B$, $Y=\Sigma^2B$, and
$C_mX=\Omega^m\Sigma^mX$.  Let $u_4:S^4\to Y$ be the
bottom map.  The generalized Whitehead product defines an actual map
\[
 h_0=[u_4,\id_Y]:\Sigma^5B=\Sigma^4X\longrightarrow Y.
 \tag{D.2.29a}
\]
Its double adjunct sends $\sigma^3b$, where $b\in H_3B$,
to $k=\{a,b\}\in H_6(\Omega^2Y;\FF_2)$, with no linear
$p_3$ component.  In the cofiber of $h_0$, the unique extensions
of the old degree-four class and its first square have nonzero product.
Thus the potentially dangerous fourth-suspension pattern is realized
on the \emph{entire unit}, not only on a finite skeleton.

The iterated counits give compatible canonical extensions
\[
 h_m:\Sigma^4C_mX\longrightarrow Y\qquad(0\le m\le4).
 \tag{D.2.29b}
\]
The specified canonical $h_4$ has no ordinary pointed extension
to $\Sigma^4C_6X$, and therefore none to $\Sigma^4QX$.
This conclusion fixes the restriction on $\Sigma^4C_4X$.
It does \emph{not} exclude an arbitrary ordinary extension to
$\Sigma^4QX$ that agrees only with $h_0$ on $\Sigma^4X$.
\end{claim}
\begin{proof}
All homology in the proof has coefficients $\FF_2$, and
$\beta=\operatorname{Sq}_*^1$.  Set $T=\Omega^2Y$ and
use the labels $a,b$ of degrees two and three, respectively,
with $\beta b=a$.  Write $B_0=\sigma b\in H_4X$.

\medskip\noindent
\textbf{The actual unit map and its cone.}
The source of the generalized Whitehead product of
$S^4=\Sigma S^3\to Y$ and $Y=\Sigma(\Sigma B)\to Y$ is
\[
 \Sigma(S^3\wedge\Sigma B)=\Sigma^5B.
\]
This gives \textup{(D.2.29a)} with the stated suspension indices.
The construction and its identification with the adjoint commutator
are \cite[Definitions~2.2--2.3 and Theorem~2.4,
pp.~10--13]{EH5Arkowitz1962}.  Its double adjunct is the
based generalized Browder construction on the bottom map
$S^2\to B\to T$ and the unit $B\to T$.
The mapping-cylinder construction in the remark following
Definition~2.3 applies to arbitrary based CW models, so the
infinite source creates no finite-polyhedron restriction.
For the primitive bottom classes $a,b$, its homology values are
\[
 (h_0^{\flat\flat})_*(\sigma^3a)=\{a,a\}=0,
 \qquad
 (h_0^{\flat\flat})_*(\sigma^3b)=\{a,b\}=k.
 \tag{D.2.29c}
\]
These are the geometric Browder operation and its suspension
compatibility in \cite[Chapter~III, Theorems~1.2--1.4,
pp.~216--219]{I_CohenLadaMay1976}.
The formula here is intentionally restricted to primitive inputs.
For nonprimitive labels the based family has convolution corrections;
no formula $x\mapsto\{a,x\}$ on every label is used.
On the bottom sphere the map is the ordinary Whitehead square
$[u_4,u_4]=-2\alpha$, with the orientation convention of D.2.17.
Its zero mod-two Hurewicz value in \textup{(D.2.29c)} does
not assert that this bottom homotopy element is zero.

Under one homology suspension, $k$ maps to the tensor
commutator $[\bar a,\bar b]\in H_7\Omega Y$; under the next,
that commutator maps to zero.  The group
$H_8(\Sigma^5B)$ is one-dimensional on $\sigma^5b$.
Thus the degree-eight attaching homology map is zero, whereas
the exact suspended-source cone criterion of D.2.26 evaluates
the old cup product to one on the connecting class over
$\sigma^5b$.  The source is six-connected, so the old
degree-four and degree-five cohomology extensions are unique.
Their product is therefore nonzero, with no choice-of-lift
qualification.  The absence of a $p_3$ component follows
directly from the second equality in \textup{(D.2.29c)}.

\medskip\noindent
\textbf{The compatible canonical stages.}
For $0\le m\le4$, let $\varepsilon^{(m)}$ be the counit
of $\Sigma^m\dashv\Omega^m$, evaluated at $\Sigma^mX$, and set
\[
 r_m=\Sigma^{4-m}\varepsilon^{(m)}:
       \Sigma^4C_mX\longrightarrow\Sigma^4X,
 \qquad h_m=h_0r_m.
 \tag{D.2.29d}
\]
The triangle identities give extension over the unit, and
naturality of the iterated counits gives compatibility as $m$
increases.  The fourfold adjunct of $h_4$ is exactly the
fourfold loop map
\[
 g_4=\Omega^4h_0:C_4X\longrightarrow\Omega^4Y.
 \tag{D.2.29e}
\]
In particular it preserves the $E_4$ homology operations.
If $q=(g_4)_*(B_0)\in H_4\Omega^4Y$, the twofold
homology suspension satisfies
\[
 \sigma^2q=k\in H_6T,
 \tag{D.2.29f}
\]
by adjunction and \textup{(D.2.29c)}.

\medskip\noindent
\textbf{A Bockstein class in the actual target.}
The element
\[
 \xi_1k=Q^7k\in H_{13}T
 \quad\hbox{satisfies}\quad
 \xi_1k\notin\beta H_{14}T.
 \tag{D.2.29g}
\]
Here is a complete basis argument for the nonboundary assertion.
The free $E_2$ homology theorem gives a polynomial basis whose
indecomposables are the ordinary Hall Lie words in the labels,
and their iterated outer restricted squares $\xi_1$;
see \cite[Chapter~III, Theorems~3.1--3.2 and Lemma~3.8]{I_CohenLadaMay1976}.
The Hall word $k=\{a,b\}$ is nonzero, and its outer square
$\xi_1k$ is an independent restricted-square generator.
Every outer $\xi_1$ output has odd degree $2|z|+1$.
There is consequently no such generator in degree fourteen.

A degree-fourteen polynomial monomial with at least two positive
factors has decomposable Bockstein, by Cartan.  A label has a
linear combination of labels as Bockstein.  Finally, expanding
the Bockstein of an ordinary Hall bracket retains an outer
bracket.  If a restricted square occurs inside such a bracket,
the restricted-Lie identity
$\{v,\xi_1w\}=\{\{v,w\},w\}$ reduces it to ordinary
Hall words.  The top Bockstein formula
\[
 \beta\xi_1z=(|z|\bmod2)z^2+\{\beta z,z\}
 \tag{D.2.29h}
\]
and the bracket Cartan formula introduce only ordinary brackets
and products in this expansion.  No term can become the
independent \emph{outer} restricted square $\xi_1k$.
These exhaust the polynomial basis in degree fourteen and prove
\textup{(D.2.29g)}.  The identities used here are
\cite[Chapter~III, Theorems~1.2--1.4,
pp.~216--219]{I_CohenLadaMay1976}.

\medskip\noindent
\textbf{The sixth-stage contradiction.}
Suppose an ordinary pointed map
$G:C_6X\to\Omega^4Y$ restricts to $g_4$ on $C_4X$.
The class $Q^7B_0\in H_{11}C_4X$ is the top $E_4$
operation on $B_0$.  In $C_6X$, $Q^8B_0$ is below the
top operation, so the ordinary Nishida Bockstein formula gives
\[
 \beta Q^8B_0=Q^7B_0.
 \tag{D.2.29i}
\]
By preservation of operations for the specified $E_4$ map
$g_4$, and homology-suspension compatibility,
\[
 \begin{split}
 \sigma^2G_*(Q^7B_0)
   &=\sigma^2(g_4)_*(Q^7B_0)\\
   &=\sigma^2Q^7q
     =Q^7(\sigma^2q)=Q^7k=\xi_1k.
 \end{split}
 \tag{D.2.29j}
\]
Naturality of the Bockstein for the \emph{ordinary} map $G$
and for homology suspension, together with
\textup{(D.2.29i)}, makes the same class equal to
\[
 \beta\bigl(\sigma^2G_*(Q^8B_0)\bigr),
\]
contrary to \textup{(D.2.29g)}.
No preservation of multiplication or Dyer--Lashof operations
has been assumed for $G$.

The use of $C_6X$ is deliberate.  At $C_5X$, $Q^8B_0$
is the top operation; its Bockstein formula has a Browder
correction involving $\beta B_0$ and $B_0$.
The argument does not discard that correction and makes no
nonextension assertion at the fifth stage.
Finally, if a map agrees only on the unit, its restriction to
$C_4X$ need not be the $E_4$ map $g_4$.
Equation \textup{(D.2.29j)} is then unavailable.
The asserted obstruction therefore has exactly the fixed
canonical scope stated in the claim.
\end{proof}
